% ==================================================================
% id:       app-011
% topic:    שימושי הנגזרת
% subtopic: משפטי ערך ממוצע (רול, לגרנז׳, קושי)
% solution: full
% ==================================================================
\textbf{איזה מההיגדים הבאים בהכרח נכון? מצאו דוגמה נגדית לכל היגד שגוי.}
\bE
 \item אם $f$ רציפה בקטע $(a,b)$ וגזירה בו, וכן $f(a)=f(b)$ אז קיימת $c\in(a,b)$ כך ש $f'(c)=0$.
\hint{השוו את הנתונים לתנאי משפט רול. אם נראה לכם שהטענה לא נכונה, חפשו דוגמה נגדית פשוטה; אחרת נסו להוכיח אותה.}
\begin{solbox}
הטענה אינה נכונה.
אילו
$f$
הייתה רציפה בקטע
$[a,b]$
היינו מיישמים את משפט רול והטענה הייתה נכונה. ננסה למצוא דוגמא לפונקציה שלא רציפה לפחות באחד הקצוות.
להלן דוגמא נגדית
$$.f(x) =\case{cc}{x, & 0<x \leq 1 \\ 1, & x = 0}$$
בבירור
$f$
גזירה בקטע
$(0,1)$
(ובפרט, רציפה בו).
מתקיים
$$,f(0)=1=f(1)$$
אבל לכל
$0<x<1$
מתקיים
$f'(x)=1\neq 0$.
\end{solbox}
\item אם $f$ גזירה בקטע $(a,b)$ ולא רציפה ב-$[a,b]$ וכן $f(a)=f(b)$ אז לכל $c\in(a,b)$ מתקיים $f'(c)\neq 0$.
\hint{השוו את הנתונים לתנאי משפט רול, וזכרו שהמשפט נותן תנאי מספיק בלבד. אם נראה לכם שהטענה לא נכונה, חפשו דוגמה נגדית פשוטה; אחרת נסו להוכיח אותה.}
\begin{solbox}
הטענה אינה נכונה. להלן דוגמא נגדית
$$.f(x) =\case{cc}{x^2, & -1<x \leq 2 \\ 4, & x = -1}$$
בבירור
$f$
לא רציפה בקטע
$[-1,2]$
(באילו נקודות היא לא רציפה?)
מתקיים
$$,f(-1)=4=f(2)$$
ו-
$f'(0)=0$.
\end{solbox}

\item אם $f$ רציפה ב $[a,b]$ וגזירה ב $(a,b)$ ומתקיים $f(a)=f(b)=0$ אז יש $c\in(a.b)$ כך ש $c^2\cdot f'(c)+2c\cdot f(c)=0$.
\hint{נסו לזהות את $x^2f'(x)+2xf(x)$ כנגזרת של פונקציה, ובדקו אם פונקציה זו מקיימת את תנאי משפט רול. אם נראה לכם שהטענה לא נכונה, חפשו דוגמה נגדית פשוטה.}
\begin{solbox}
הטענה נכונה.
נשים לב כי
$$,2xf(x)+x^2f'(x)=\left(x^2f(x)\right)'$$
לכן
נתבונן בפונקציה
$$g(x)=x^2f(x)$$
וננסה למצוא
$c$
כלשהו שמאפס את נגזרתה.
זוהי פונקציה רציפה בקטע
$[a,b]$
וגזירה בקטע
$(a,b)$
כמכפלה של פונקציות כאלה.
בנוסף
$$.g(a)=a^2f(a)=0=b^2f(b)=g(b)$$
ממשפט רול, קיים
$c\in(a,b)$
כך ש-
$$,0=g'(c)=2cf(c)+c^2f'(c)$$
כנדרש.
\end{solbox}
\item אם $f$ רציפה ב $[a,b]$ וגזירה ב $(a,b)$ ומתקיים $f(a)=f(b)$ אז קיימת $c\in(a,b)$ כך ש $e^c\cdot f(c)+e^c\cdot f'(c)=0$.
\hint{נסו לזהות את $e^xf(x)+e^xf'(x)$ כנגזרת של פונקציה, ובדקו אם פונקציה זו מקיימת את תנאי משפט רול. אם נראה לכם שהטענה לא נכונה, חפשו דוגמה נגדית פשוטה.}
\begin{solbox}
הטענה אינה נכונה.
 ניקח
 $f(x)=1$
 (פונקציה קבועה)
 בקטע
 $[a,b]$.
 בבירור
 $f$
 עונה על התנאים של השאלה ומתקיים
 $f'(x)=0$
 לכל
 $x\in[a,b]$,
 בפרט
 $f(x)+f'(x)=1\neq0$
 לכל
 $x\in[a,b]$,
 לכן גם
 $e^xf(x)+e^xf'(x)=e^x\neq0$
 לכל
 $x\in[a,b]$.

 מה היה משתבש בפתרון אם היינו מגדירים פונקציה
 $g(x)=e^xf(x)$
 והיינו מנסים להוכיח בדומה למה שהוכחנו?
\end{solbox}
\eE
